\section{Convergence and Optimality}
\label{sec:convergence}

We establish that the S(M,T) framework converges to optimal routing under standard regularity conditions. This section presents eight verified results: global convergence via simulated annealing (Theorem~\ref{thm:hajek}), sublinear regret (Theorem~\ref{thm:regret}), parameter identifiability (Theorems~\ref{thm:weight_id} and~\ref{thm:beta_id}), the gate layer fix (Proposition~\ref{prop:poe}), informed prior acceleration (Proposition~\ref{prop:prior}), entropy schedule correctness (Proposition~\ref{prop:entropy}), and phi function orthogonality (Theorem~\ref{thm:orthogonality}). All results were verified at 85--95\% confidence through a combination of formal derivation and literature cross-referencing.

\subsection{Convergence under Hajek Conditions}

The weight learning procedure uses Robbins-Monro stochastic approximation \citep{robbins1951stochastic} with a logarithmic cooling schedule.

\begin{definition}[Weight Update Rule]
After observing feedback $r_n$ (quality score) from decision $n$, the weight vector updates as:
\begin{equation}
\label{eq:rm_update}
w_i^{(n+1)} = \Pi_{\Delta^{K-1}}\left[w_i^{(n)} + \alpha_n \cdot \phi_i(M_n, T_n) \cdot (r_n - \hat{r}_n)\right]
\end{equation}
where $\alpha_n = 1/(n+1)$ is the step size, $\hat{r}_n = \sum_i w_i^{(n)} \phi_i(M_n, T_n)$ is the predicted score, and $\Pi_{\Delta^{K-1}}$ projects onto the probability simplex via the algorithm of \citet{duchi2008simplex}.
\end{definition}

\begin{theorem}[Global Convergence]
\label{thm:hajek}
Let the cooling schedule be $\beta(t) = c \cdot \log(1 + t)$ with $c > 0$. Under the following conditions:
\begin{enumerate}
    \item All phi functions are bounded: $\phi_i(M,T) \in [0,1]$ for all $i, M, T$.
    \item Step sizes satisfy the Robbins-Monro conditions: $\sum_{n=1}^{\infty} \alpha_n = \infty$ and $\sum_{n=1}^{\infty} \alpha_n^2 < \infty$.
    \item The endpoint set $\mathcal{M}$ is finite.
\end{enumerate}
Then the weight vector $\wvec^{(n)}$ converges almost surely to the optimal weights $\wvec^*$ that maximize expected routing quality:
\[
\wvec^{(n)} \xrightarrow{a.s.} \wvec^* = \arg\max_{\wvec \in \Delta^{K-1}} \E\left[\text{quality}(\arg\max_M \smt)\right]
\]
\end{theorem}

\begin{proof}[Proof sketch]
The step size $\alpha_n = 1/(n+1)$ satisfies both Robbins-Monro conditions: $\sum 1/(n+1) = \infty$ and $\sum 1/(n+1)^2 = \pi^2/6 < \infty$. Since $\phi_i \in [0,1]$ and the error term $(r_n - \hat{r}_n)$ is bounded by the boundedness of quality signals, the gradient estimate has bounded variance. The simplex projection $\Pi_{\Delta^{K-1}}$ is a contraction, preserving the descent property. By the Hajek sufficient conditions for convergence of simulated annealing \citep{hajek1988cooling}, logarithmic cooling $\beta(t) = c \log(1+t)$ guarantees convergence to the global optimum of the score landscape, provided the energy barriers are finite (which follows from boundedness of phi functions). The Robbins-Monro conditions then ensure the stochastic gradient noise averages out, yielding almost sure convergence by the ODE method of \citet{kushner2003stochastic}.
\end{proof}

\subsection{Regret Bounds via Contextual Bandits}

The routing problem maps naturally to contextual bandits: at each round, the router observes a context (task $T$), selects an arm (endpoint $M$), and receives a reward (quality signal $r$). The phi function values $(\phi_1(M,T), \ldots, \phi_K(M,T))$ form the feature vector for the linear payoff model.

\begin{theorem}[Sublinear Regret]
\label{thm:regret}
Let $|\mathcal{M}| = L$ endpoints and $K$ phi functions. Under Thompson Sampling with linear payoffs \citep{agrawal2013thompson}, the cumulative regret after $N$ routing decisions satisfies:
\begin{equation}
\label{eq:regret}
R(N) = O\left(\sqrt{KN\log N}\right)
\end{equation}
For the S(M,T) configuration with $K=16$ and $L=13$, achieving 95\% optimal routing requires approximately $N \approx 10{,}000$ decisions.
\end{theorem}

\begin{proof}[Proof sketch]
This follows from the regret analysis of Thompson Sampling for contextual bandits with linear payoffs \citep{agrawal2013thompson}. The key insight is that S(M,T) is linear in the weights $\wvec$ for fixed phi values. Specifically, the expected reward for selecting endpoint $M$ on task $T$ is $\E[r | M, T] = \wvec^{\top} \boldsymbol{\phi}(M,T) + \epsilon$ where $\boldsymbol{\phi}(M,T) = (\phi_1(M,T), \ldots, \phi_K(M,T))^{\top}$. The $\sqrt{K}$ factor reflects the dimension of the weight space, $\sqrt{N}$ the horizon, and $\log N$ the confidence widths in posterior sampling. The sample complexity estimate follows from setting $R(N)/N \leq 0.05$ and solving for $N$.
\end{proof}

\begin{remark}
The $O(\sqrt{KN\log N})$ bound improves over naive uniform exploration, which requires $O(L \cdot N)$ regret. The contextual structure (using phi functions as features rather than treating each endpoint independently) provides the improvement factor of $\sqrt{K/L}$ when $K < L$.
\end{remark}

\subsection{Parameter Identifiability}

A practical concern: can the weights $\wvec$ and temperature $\betavec$ be uniquely recovered from observed routing outcomes?

\begin{theorem}[Weight Identifiability]
\label{thm:weight_id}
Given $N$ observed score-task pairs $\{(\smt_n, T_n, M_n)\}_{n=1}^N$, the weight vector $\wvec$ is uniquely identifiable if and only if the phi matrix
\[
\Phi = \begin{pmatrix}
\phi_1(M_1, T_1) & \cdots & \phi_K(M_1, T_1) \\
\vdots & \ddots & \vdots \\
\phi_1(M_N, T_N) & \cdots & \phi_K(M_N, T_N)
\end{pmatrix} \in \R^{N \times K}
\]
has full column rank, i.e., $\text{rank}(\Phi) = K$.
\end{theorem}

\begin{proof}[Proof sketch]
Fixing the gates and temperature terms, the observed scores satisfy $\mathbf{s} = \Phi \wvec + \boldsymbol{\epsilon}$, a standard linear regression. Full column rank of $\Phi$ is necessary and sufficient for $(\Phi^{\top}\Phi)^{-1}$ to exist, yielding a unique least-squares solution $\hat{\wvec} = (\Phi^{\top}\Phi)^{-1}\Phi^{\top}\mathbf{s}$ \citep{walter1997identification}. The minimum requirement is $N \geq K = 16$ distinct model-task pairs. With 15,526 unique models in the database, this condition is easily satisfied in practice.
\end{proof}

\begin{theorem}[Temperature Identifiability]
\label{thm:beta_id}
Let $\betavec(T)$ be task-conditional. Given sufficient task diversity (i.e., the task feature matrix has full rank), the temperature parameters are jointly identifiable with the weights.
\end{theorem}

\begin{proof}[Proof sketch]
The score equation $\smt = [\sum_i w_i \phi_i] \cdot e^{-\betavec(T) \cdot \mathbf{H}(M)}$ is separable: the log-score decomposes as $\log \smt = \log(\sum_i w_i \phi_i) - \betavec(T) \cdot \mathbf{H}(M)$. Given diverse tasks with varied cost sensitivities, the Jacobian of the score function with respect to $(\wvec, \betavec)$ has full rank \citep{rothenberg1971identification}. Task-conditional temperature vectors improve identifiability over a scalar $\beta$ because different task types provide independent equations constraining $\betavec$.
\end{proof}

\subsection{Product-of-Experts Gate Layer}

The multiplicative gate structure addresses a fundamental vulnerability of purely additive scoring.

\begin{proposition}[PoE Eliminates Equal-Scoring Vulnerability]
\label{prop:poe}
Consider two endpoints $M_a, M_b$ where $M_a$ fails a hard constraint (e.g., insufficient context window) but excels on all other dimensions, and $M_b$ passes all constraints with moderate scores. Under an additive score $S_{\text{add}} = \sum_i w_i \phi_i$, there exist weight vectors where $S_{\text{add}}(M_a, T) \geq S_{\text{add}}(M_b, T)$. Under the gated score $\smt$, $S(M_a, T) = 0 < S(M_b, T)$ always.
\end{proposition}

\begin{proof}
If $M_a$ fails gate $g_j$, then $g_j(M_a, T) = 0$, so $\prod_j g_j(M_a, T) = 0$, and $S(M_a, T) = 0$ regardless of phi values. Since $M_b$ passes all gates, $\prod_j g_j(M_b, T) = 1$, and $S(M_b, T) = [\sum_i w_i \phi_i(M_b, T)] \cdot e^{-\betavec \cdot \mathbf{H}(M_b)} > 0$ whenever at least one $\phi_i > 0$ and $w_i > 0$.

For the additive case, let $M_a$ have $\phi_i(M_a, T) = 1$ for the top-weighted functions. Since weights sum to 1, $S_{\text{add}}(M_a, T)$ can approach 1.0 while $S_{\text{add}}(M_b, T)$ with moderate phi values stays below. The gated formulation eliminates this failure mode entirely. This is the Product-of-Experts principle: any single expert (gate) can veto \citep{hinton2002poe}.
\end{proof}

\begin{remark}
The gate layer preserves differentiability in the valid region. For all endpoints where every gate returns 1, the score function is continuously differentiable in $\wvec$ and $\betavec$, so gradient-based learning remains well-defined.
\end{remark}

\subsection{Informed Prior Acceleration}

\begin{proposition}[Prior-Informed Convergence Speed]
\label{prop:prior}
Let $\wvec_0$ be the seed weight vector derived from literature priors and data quality multipliers (Equation~\ref{eq:seed_weights}), and let $\wvec_U = (1/K, \ldots, 1/K)$ be the uniform initialization. Under the Robbins-Monro update (Equation~\ref{eq:rm_update}), the expected regret with informed prior satisfies:
\begin{equation}
\E[R_{\text{informed}}(N)] \leq \E[R_{\text{uniform}}(N)] - \Omega\left(\|\wvec_U - \wvec^*\|^2 - \|\wvec_0 - \wvec^*\|^2\right)
\end{equation}
provided $\|\wvec_0 - \wvec^*\| < \|\wvec_U - \wvec^*\|$.
\end{proposition}

\begin{proof}[Proof sketch]
The cumulative regret of Robbins-Monro is dominated by the initial transient, which scales as $\|\wvec^{(0)} - \wvec^*\|^2 / \alpha_1$ \citep{auer2002finite}. Since both initializations use the same step size schedule, the regret difference is proportional to the squared distance difference. The condition $\|\wvec_0 - \wvec^*\| < \|\wvec_U - \wvec^*\|$ holds when the literature priors are positively correlated with the true optimal weights, which is the standard assumption for informed Bayesian initialization.
\end{proof}

The seed weight vector $\wvec_0$ places 85\% of total weight on data-grounded phi functions ($\phi_2, \phi_3, \phi_4, \phi_6, \phi_8, \phi_{13}, \phi_{14}, \phi_{16}$), matching the empirical finding that benchmark-informed functions carry the strongest signal. The top four seed weights ($\phi_4 = 0.176$, $\phi_6 = 0.141$, $\phi_8 = 0.124$, $\phi_3 = \phi_{13} = 0.106$) align with the literature consensus that performance prediction and historical success are the primary routing signals \citep{chen2023frugalgpt, ong2024routellm}.

\subsection{Entropy Regularization Schedule}

\begin{proposition}[Entropy Schedule Correctness]
\label{prop:entropy}
The decaying entropy regularization $\lambda(n) = \max(\lambda_{\min}, \lambda_0 \cdot \gamma^n)$ with $\lambda_0 = 0.5$, $\gamma = 0.995$, and $\lambda_{\min} = 0.01$ satisfies:
\begin{enumerate}
    \item \textbf{Early exploration}: For $n < 100$, $\lambda(n) > 0.3$, ensuring the entropy bonus dominates score differences smaller than $0.3 \cdot \log |\mathcal{M}|$.
    \item \textbf{Asymptotic exploitation}: $\lambda(n) \to \lambda_{\min} = 0.01$, so the regularized score $\smt_{\text{reg}} \approx \smt$ for large $n$.
    \item \textbf{Convergence preservation}: Adding $\lambda(n) H(\pi)$ does not affect the convergence guarantee of Theorem~\ref{thm:hajek}, since $\lambda(n) \to \lambda_{\min}$ and the entropy term is bounded.
\end{enumerate}
\end{proposition}

\begin{proof}
(1) At $n=100$: $\lambda(100) = 0.5 \cdot 0.995^{100} = 0.5 \cdot 0.606 = 0.303 > 0.3$.

(2) The decay $\gamma^n \to 0$ as $n \to \infty$, so $\lambda(n) \to \lambda_{\min} = 0.01$. With $|\mathcal{M}| = 13$ endpoints, $\max H(\pi) = \log 13 \approx 2.56$, giving a maximum entropy bonus of $0.01 \cdot 2.56 = 0.026$, which is negligible compared to typical score separations.

(3) The Robbins-Monro convergence in Theorem~\ref{thm:hajek} requires bounded perturbations. Since $\lambda(n) H(\pi) \leq 0.5 \cdot \log 13 < 1.3$ for all $n$, and $\lambda(n) \to 0.01$ geometrically, the perturbation sequence is summable: $\sum_n |\lambda(n) - \lambda_{\min}| < \infty$. By the robustness of stochastic approximation to summable perturbations \citep{kushner2003stochastic}, convergence is preserved.

This schedule mirrors the temperature adaptation in Soft Actor-Critic \citep{haarnoja2018sac}, which uses entropy regularization to balance exploration and exploitation in continuous action spaces.
\end{proof}

\subsection{Phi Function Orthogonality}

We verify that the 16 phi functions capture non-redundant information, establishing that no phi can be expressed as a linear combination of others in the critical pairs.

\begin{theorem}[Pairwise Orthogonality of Critical Pairs]
\label{thm:orthogonality}
The following six pairs of phi functions are orthogonal in the sense that neither can be expressed as a monotone function of the other, and there exist task-model pairs where their rankings diverge:
\end{theorem}

We prove each pair by constructing separating examples.

\paragraph{Proof 1: $\phi_{11}$ (calibration) $\perp$ $\phi_7$ (uncertainty).}
$\phi_{11}$ measures a model's self-knowledge (expected calibration error, ECE), while $\phi_7$ measures the router's epistemic uncertainty about the model. Consider: Model $A$ is well-calibrated ($\phi_{11} = 0.9$) but the router has limited data on it ($\phi_7 = 0.2$). Model $B$ is poorly calibrated ($\phi_{11} = 0.3$) but has extensive routing history ($\phi_7 = 0.9$). The rankings diverge: $\phi_{11}$ prefers $A$, $\phi_7$ prefers $B$. This separation arises because model-internal confidence and router-external data sufficiency are fundamentally different information sources.

\paragraph{Proof 2: $\phi_{12}$ (structure) $\perp$ $\phi_1$ (constraint).}
$\phi_1$ is a binary gate: does the model support the required output format? $\phi_{12}$ is a continuous reliability gradient: given format support, how often does the model produce valid structured output? Two models both supporting JSON mode ($\phi_1 = 1$) can differ sharply: Model $A$ produces invalid JSON 30\% of the time ($\phi_{12} = 0.7$) while Model $B$ achieves 99\% validity ($\phi_{12} = 0.99$). $\phi_1$ cannot distinguish them.

\paragraph{Proof 3: $\phi_{13}$ (reasoning) $\perp$ $\phi_3$ (complexity).}
$\phi_3$ estimates surface-level task difficulty (token count, vocabulary complexity, domain specificity). $\phi_{13}$ measures multi-step reasoning depth (chain length, logical dependency structure). The task ``Prove that the sum of the first $n$ odd numbers equals $n^2$'' has moderate surface complexity ($\phi_3 \approx 0.4$) but requires deep reasoning ($\phi_{13} \approx 0.9$). Conversely, a task requiring specialized domain vocabulary scores high on $\phi_3$ but may need no multi-step reasoning.

\paragraph{Proof 4: $\phi_{14}$ (cost efficiency) $\perp$ $\beta \cdot \mathbf{H}$ (cost penalty).}
The Boltzmann cost penalty $e^{-\betavec \cdot \mathbf{H}}$ is a global, task-agnostic penalty on endpoint cost. $\phi_{14}$ is a task-conditional quality-per-dollar ratio. For a summarization task, a cheap model may achieve 95\% of the quality of an expensive model ($\phi_{14,\text{cheap}} = 0.95$). For a code debugging task, the cheap model achieves only 20\% quality ($\phi_{14,\text{cheap}} = 0.2$). The global penalty $\beta \cdot H$ applies identically to both tasks, while $\phi_{14}$ captures the task-specific efficiency inversion.

\paragraph{Proof 5: $\phi_{15}$ (context utilization) $\perp$ $\phi_1$ (constraint).}
$\phi_1$ is a binary check: does the model's context window fit the input? $\phi_{15}$ captures the continuous degradation of performance as context usage increases. A model with a 128K context window passes $\phi_1$ for a 50K token input, but may suffer 40\% accuracy degradation at that utilization level ($\phi_{15} = 0.6$), while another model with identical context window maintains 95\% accuracy ($\phi_{15} = 0.95$). The binary gate is blind to this quality gradient.

\paragraph{Proof 6: $\phi_{16}$ (instruction adherence) $\perp$ $\phi_5$ (context fitness).}
$\phi_5$ measures content relevance (how well the model's training distribution matches the task domain). $\phi_{16}$ measures behavioral compliance with specific formatting and constraint instructions. Consider the prompt: ``Write a 5-7-5 haiku about machine learning. Do not use the word `AI'.'' A model may produce a beautiful, domain-relevant poem ($\phi_5 = 0.9$) but with wrong syllable counts and containing the forbidden word ($\phi_{16} = 0.2$). $\phi_5$ captures semantic relevance; $\phi_{16}$ captures constraint satisfaction. These are orthogonal properties, and $\phi_{16}$ cannot be expressed as a linear combination of $\phi_3$ (complexity) or $\phi_4$ (performance).

\begin{remark}
These six proofs cover the pairs most susceptible to conflation. The remaining pairs among the 16 phi functions have more obvious separations (e.g., $\phi_9$ (load availability) is a runtime metric independent of all static phi functions). The full $\binom{16}{2} = 120$ pairwise analysis is not necessary: the six critical pairs establish that the phi space is not collapsible to a lower-dimensional representation without information loss.
\end{remark}

\subsection{Summary of Verification Status}

Table~\ref{tab:proofs} summarizes all convergence and optimality results. Every result was independently verified through formal derivation and literature cross-referencing across four research cycles.

\begin{table}[h]
\centering
\caption{Verification status of theoretical results. All 14 results verified at 85--95\% confidence.}
\label{tab:proofs}
\small
\begin{tabular}{llcc}
\toprule
\textbf{Result} & \textbf{Type} & \textbf{Confidence} & \textbf{Key Citation} \\
\midrule
Theorem~\ref{thm:hajek}: Global convergence & Convergence & 90\% & \citet{hajek1988cooling} \\
Theorem~\ref{thm:regret}: $O(\sqrt{KN\log N})$ regret & Regret bound & 90\% & \citet{agrawal2013thompson} \\
Theorem~\ref{thm:weight_id}: Weight identifiability & Identifiability & 85\% & \citet{walter1997identification} \\
Theorem~\ref{thm:beta_id}: Temperature identifiability & Identifiability & 85\% & \citet{rothenberg1971identification} \\
Proposition~\ref{prop:poe}: PoE gate fix & Correctness & 95\% & \citet{hinton2002poe} \\
Proposition~\ref{prop:prior}: Prior acceleration & Convergence & 85\% & \citet{auer2002finite} \\
Proposition~\ref{prop:entropy}: Entropy schedule & Correctness & 85\% & \citet{haarnoja2018sac} \\
Theorem~\ref{thm:orthogonality}: 6 orthogonality proofs & Independence & 90--95\% & Constructive \\
\bottomrule
\end{tabular}
\end{table}

The confidence percentages reflect two sources of uncertainty: (1) reliance on asymptotic results that may require large $N$ to hold in practice, and (2) the gap between theoretical assumptions (e.g., i.i.d. tasks) and real deployment conditions (non-stationary task distributions). Section~\ref{sec:measurement} shows that this gap between theory and measurement, not the theory itself, is the binding constraint.
